% HQIV vs TUFT shell language (include in TUFT / hadron manuscripts)
%
% Usage: % HQIV vs TUFT shell language (include in TUFT / hadron manuscripts)
%
% Usage: % HQIV vs TUFT shell language (include in TUFT / hadron manuscripts)
%
% Usage: % HQIV vs TUFT shell language (include in TUFT / hadron manuscripts)
%
% Usage: \input{../include/shell_ontology_messaging.tex}
% (adjust path relative to your paper directory)

\paragraph{Shell index ontology (do not conflate).}
HQIV names \texttt{referenceM} for the substrate lock-in shell used in cosmology,
CMB readouts, and the single-scale mass witness.
TUFT names \emph{Hopf windings} ($n=1,2,3$ on the weak, strong, and heavy sectors)
and the corresponding \emph{Beltrami chart rows} $m=n+1$ (today $m=2,3,4$).
Baryon spectroscopy is indexed on the \textbf{heavy TUFT chart}
$m_{\rm hadron}(n,\ell)=m_{\rm heavy}+n+\ell$ with internal quanta $(n,\ell)$;
these are not the same labels as HQIV lock-in bookkeeping even when
$m_{\rm heavy}=\mathrm{referenceM}$ numerically at the current pins.
Meson vector grounds use the strong/heavy chart ratio $3/4$.
When comparing to data, state which chart is primary; do not write
``shell $m=\mathrm{referenceM}$'' alone in TUFT hadron sections.

% (adjust path relative to your paper directory)

\paragraph{Shell index ontology (do not conflate).}
HQIV names \texttt{referenceM} for the substrate lock-in shell used in cosmology,
CMB readouts, and the single-scale mass witness.
TUFT names \emph{Hopf windings} ($n=1,2,3$ on the weak, strong, and heavy sectors)
and the corresponding \emph{Beltrami chart rows} $m=n+1$ (today $m=2,3,4$).
Baryon spectroscopy is indexed on the \textbf{heavy TUFT chart}
$m_{\rm hadron}(n,\ell)=m_{\rm heavy}+n+\ell$ with internal quanta $(n,\ell)$;
these are not the same labels as HQIV lock-in bookkeeping even when
$m_{\rm heavy}=\mathrm{referenceM}$ numerically at the current pins.
Meson vector grounds use the strong/heavy chart ratio $3/4$.
When comparing to data, state which chart is primary; do not write
``shell $m=\mathrm{referenceM}$'' alone in TUFT hadron sections.

% (adjust path relative to your paper directory)

\paragraph{Shell index ontology (do not conflate).}
HQIV names \texttt{referenceM} for the substrate lock-in shell used in cosmology,
CMB readouts, and the single-scale mass witness.
TUFT names \emph{Hopf windings} ($n=1,2,3$ on the weak, strong, and heavy sectors)
and the corresponding \emph{Beltrami chart rows} $m=n+1$ (today $m=2,3,4$).
Baryon spectroscopy is indexed on the \textbf{heavy TUFT chart}
$m_{\rm hadron}(n,\ell)=m_{\rm heavy}+n+\ell$ with internal quanta $(n,\ell)$;
these are not the same labels as HQIV lock-in bookkeeping even when
$m_{\rm heavy}=\mathrm{referenceM}$ numerically at the current pins.
Meson vector grounds use the strong/heavy chart ratio $3/4$.
When comparing to data, state which chart is primary; do not write
``shell $m=\mathrm{referenceM}$'' alone in TUFT hadron sections.

% (adjust path relative to your paper directory)

\paragraph{Shell index ontology (do not conflate).}
HQIV names \texttt{referenceM} for the substrate lock-in shell used in cosmology,
CMB readouts, and the single-scale mass witness.
TUFT names \emph{Hopf windings} ($n=1,2,3$ on the weak, strong, and heavy sectors)
and the corresponding \emph{Beltrami chart rows} $m=n+1$ (today $m=2,3,4$).
Baryon spectroscopy is indexed on the \textbf{heavy TUFT chart}
$m_{\rm hadron}(n,\ell)=m_{\rm heavy}+n+\ell$ with internal quanta $(n,\ell)$;
these are not the same labels as HQIV lock-in bookkeeping even when
$m_{\rm heavy}=\mathrm{referenceM}$ numerically at the current pins.
Meson vector grounds use the strong/heavy chart ratio $3/4$.
When comparing to data, state which chart is primary; do not write
``shell $m=\mathrm{referenceM}$'' alone in TUFT hadron sections.
